% =====================================================================
%  Da incollare in Overleaf IMMEDIATAMENTE PRIMA di \chapter{Conclusions}.
%  LaTeX rinumera da solo: questo diventa il capitolo 12 e Conclusions il 13.
%  Le tabelle usano booktabs (\toprule/\midrule/\bottomrule).
%
%  RIFERIMENTI A PARTI ESISTENTI DELLA TESI
%  Questo capitolo cita cinque punti del testo gia' scritto. Se quelle parti
%  non hanno un \label, LaTeX stampera' "??": in tal caso o aggiungi il \label
%  indicato, oppure sostituisci il \ref con il numero letterale.
%
%    \ref{ch:results}      -> cap. 8   "Large-scale analysis results"
%    \ref{sec:manual}      -> sez. 8.1.3 "Manual audit results"
%    \ref{ch:discussion}   -> cap. 9   "Discussion"
%    \ref{ch:ethics}       -> cap. 11  "Ethical discussion"
%    \ref{sec:llm}         -> sez. 10.3 (lo Stage 2B su Ollama/llama3)
%                             usato solo nel file etica_sostituzione.tex
% =====================================================================

\chapter{A Longitudinal Re-Analysis of the Ecosystem}
\label{ch:longitudinal}

The measurement reported in chapter~\ref{ch:results} describes the MCP ecosystem
as it was at a single moment in time. Two months after that first execution we
ran SAMS a second time, over the same 69,104 servers and with the same
parameters, for two distinct reasons. The first is internal: repeating a
large-scale measurement is the only way to establish how much of its outcome is
a property of the ecosystem and how much is a property of the instrument. The
second is external: two snapshots of the same population, taken at a distance,
make it possible to ask questions that a single snapshot cannot answer, and in
particular whether MCP servers modify their declared behaviour after the moment
at which a user would have approved them. This chapter reports what the second
measurement found, what it says about the reliability of the first, and what a
longitudinal view adds. It also documents, in
section~\ref{sec:long:threats}, a set of measurement artefacts that we
encountered and corrected, because they bear directly on how results of this
kind should be interpreted.

\section{Design of the second measurement}
\label{sec:long:design}

The second execution reused the dataset, the framework selection, the framework
parameters and the three-stage post-processing of the first, and was distributed
over the same nine virtual machines. Holding the pipeline fixed was a deliberate
choice: any divergence in the results is then attributable either to the
ecosystem itself or to a small number of deviations that we were forced to make
and that we enumerate below. Where the two executions cannot be compared on
equal terms, we say so rather than presenting a difference that the design does
not support.

\section{Declared deviations from the original setup}
\label{sec:long:deviations}

Three aspects of the pipeline could not be preserved unchanged, and each of them
constrains the comparison in a different way.

\paragraph{One framework is no longer executable.}
\texttt{mcp-scan} delegates the analysis that produces its \texttt{E001} and
\texttt{W001} findings to a remote service operated by Invariant Labs. That
service no longer exists: the public endpoint has no DNS record, and the command
line interface falls back on a Snyk endpoint that answers every request with
\texttt{Push key is required}. The logs of the second run contain 110,658 such
failures, one for every analysed tool, and consequently every tool is reported
as safe. This is not a transient outage but the consequence of an acquisition,
and it cannot be recovered in post-processing: pinning an older release does not
help, because the server the release depended on has been decommissioned. We
therefore imported the \texttt{mcp-scan} results of the first execution into the
second, marking every such finding with an explicit provenance field. The
consequence must be kept in mind when reading the tables: for the categories
that depend on \texttt{mcp-scan} the comparison is, by construction, a
comparison of a dataset with itself.

This is also a result in its own right. A scanner that delegates its detection
logic to a proprietary backend stops working when that backend disappears, and
the reproducibility of any measurement built on it expires with the vendor's
commercial decisions.

\paragraph{The Stage 2B classifier changed.}
In the first execution the residual \textsc{uncertain} bucket was resolved by a
local Ollama instance serving \texttt{llama3}. In the second execution it was
resolved by Claude Sonnet, a hosted model, because the local path was not
available on the machine that carried out the post-processing. The change is
material and we state it explicitly, together with its consequences. The payload
sent to the hosted model consisted of the derived findings, not of the full
source of the analysed servers; those findings, however, include the evidence
field of the \texttt{credential-leak} category, and therefore third-party
credentials embedded in public repositories were transmitted to an external
provider. No credential was ever exercised against the corresponding service,
before or after classification. The ethical implications of this choice are
discussed in chapter~\ref{ch:ethics}.

The change is also less consequential for the results than it may appear. The
deterministic high-confidence rules of Stage 2A are identical in the two
executions and involve no model at all, and Stage 2B intervenes on the 6,039
uncertain findings that remain out of the 56,908 surviving Stage 1, promoting
668 of them to confirmed. The classifier therefore accounts for 2.8\% of the
24,225 confirmed findings of the second measurement, against a comparable 2.0\%
for \texttt{llama3} in the first, and the remainder is produced by rules that
did not change. Whatever the merits of the two models, the outcome of this study
does not rest on them.

\paragraph{The dataset was covered in a single pass.}
The first execution analysed the 60,205 GitHub servers and the NPX packages in
two separate runs that were merged afterwards, whereas the second covered the
69,104 servers of the unified dataset in one pass. The totals are therefore
comparable, but the per-stage intermediate counts of the first execution are
distributed differently.

\section{Results of the second measurement}
\label{sec:long:results}

After post-processing, the second execution retained 24,225 confirmed findings
over 11,487 distinct servers, against the 27,958 findings over 11,841 servers of
the first. Table~\ref{tab:long:funnel} places the two funnels side by side.

\begin{table}[htbp]
\centering
\caption{The two measurements, stage by stage}
\label{tab:long:funnel}
\begin{tabular}{lrr}
\toprule
 & First execution & Second execution \\
\midrule
Servers analysed              & 69,104      & 69,104 \\
Frameworks executed           & 7           & 6 (\texttt{mcp-scan} imported) \\
Raw findings                  & 3,102,800   & 3,531,813 \\
After Stage 1                 & 73,594      & 56,908 \\
Stage 2B (uncertain $\to$ TP) & 5,800 $\to$ 565 & 6,039 $\to$ 668 \\
\textbf{Confirmed findings}   & \textbf{27,958} & \textbf{24,225} \\
Distinct servers affected     & 11,841      & 11,487 \\
\bottomrule
\end{tabular}

\vspace{2mm}
\footnotesize
The raw count of the second execution is reported after deduplication; the
reason is explained in section~\ref{sec:long:duplicates}. The Stage 2B figure of
the second execution excludes the 1,337 uncertain findings inherited from the
imported \texttt{mcp-scan} results, which were not reclassified.
\end{table}

The aggregate stability of the two totals conceals a considerably less stable
population. Only 8,764 servers carry confirmed findings in both executions,
against 3,077 that appear only in the first and 2,723 only in the second, for a
Jaccard coefficient of 60.2\%. In other words, roughly two servers out of five
that the pipeline reports as vulnerable at any given time are not the same
servers it reported two months earlier. Part of this turnover is real change in
the ecosystem, and part of it is the volatility of which servers can be started
at all, a point we return to in section~\ref{sec:long:rugpull}.

The per-framework comparison in Table~\ref{tab:long:framework} shows that the
divergence is not evenly distributed.

\begin{table}[htbp]
\centering
\caption{Confirmed findings per framework}
\label{tab:long:framework}
\begin{tabular}{lrrr}
\toprule
Framework & First execution & Second execution & $\Delta$ \\
\midrule
\texttt{mcp-check}         & 14,983 & 10,500 & $-4,483$ \\
\texttt{mcp-guard}         &  6,010 &  7,386 & $+1,376$ \\
\texttt{mcp-scan}          &  4,396 &  4,396 & $\pm 0$ \\
\texttt{mcp-watch}         &  1,166 &  1,224 & $+58$ \\
\texttt{mcp-security-scan} &  1,374 &    651 & $-723$ \\
\texttt{mcp-shield}        &     16 &     21 & $+5$ \\
\texttt{tool\_fuzzing}     &     13 &     47 & $+34$ \\
\midrule
Total                      & 27,958 & 24,225 & $-3,733$ \\
\bottomrule
\end{tabular}

\vspace{2mm}
\footnotesize
\texttt{mcp-scan} is unchanged by construction, its findings having been
imported. The counts of the second execution are deduplicated
(section~\ref{sec:long:duplicates}); those of the first are not affected, since
its outputs contain no measurable duplication.
\end{table}

The contraction of \texttt{mcp-check} deserves a comment, because it accounts
for most of the overall decrease and could be read as an improvement in protocol
conformance. It is not. Restricting the comparison to the 23,971 GitHub servers
that both executions managed to start successfully, \texttt{mcp-check} reports
at least one finding on 19.9\% of them in the second execution against 18.2\% in
the first: on identical servers it finds more, not fewer. The failure rate
attributable to timeouts is likewise unchanged, at 89.2\% of failures against
91.3\%. What changed is which servers can be started at all, and since
\texttt{mcp-check} must complete a handshake before it can report anything, its
count follows the reachability of the ecosystem rather than its conformance.

\section{Re-assessing reliability}
\label{sec:long:reliability}

We repeated the manual audit of section~\ref{sec:manual} on the second
measurement, with the same method: a sample stratified by server across the same
seventeen categories, with each verdict decided by reading the actual source of
the server rather than the finding alone, complemented by an exhaustive audit of
the \texttt{credential-leak} class. In total 1,590 findings were inspected,
against 1,579 in the first execution. Weighting each category by its size, the
estimated precision of the second measurement is 49.7\%, against the 64.8\%
estimated for the first, which places the expected number of genuinely
exploitable findings at roughly 12,000 rather than 18,100.

A drop of fifteen points invites the conclusion that the pipeline became less
accurate, and the more interesting finding of this chapter is that it did not.
The difference is concentrated in a small number of categories and is dominated
by a single classification decision. In \texttt{credential-leak}, the largest
audited class, 54.9\% of the findings of the first execution and 53.2\% of those
of the second consist of the same pattern: the idiom
\texttt{token.write(creds.to\_json())}, produced by the official Google OAuth
quickstart, in which the user's own token is cached locally on the user's own
machine after an interactive consent flow. Counting that pattern as a false
positive puts the category at 38\%; counting it as a true positive puts it above
90\%, which reproduces the rate reported for the first execution. The audit
checklist of section~\ref{sec:manual} does not cover this case at all: it was
written for secrets committed to a repository, and is simply silent on
credentials that a program writes at runtime into its own cache.

The consequence is methodological rather than numerical. A precision figure of
this kind is not a property of the detector alone but of the detector together
with an adjudication convention, and a convention that remains implicit will
produce apparently unstable estimates across otherwise identical measurements.
We therefore report both readings explicitly, and recommend that comparable
studies state their adjudication rules for borderline classes with the same care
they devote to their detection rules.

\section{Tool-description drift and rug pulls}
\label{sec:long:rugpull}

The threat scenario TS-01 includes the case of a server that alters the
description of its tools after the user has approved it. Because the MCP consent
model operates at the granularity of the server, a description that changes
after installation is never re-approved: the agent simply reads the new text at
the next \texttt{tools/list} call. This is the class of attack that motivates
the hash store maintained by \texttt{mcp-scan}, and it is the one class that a
single measurement cannot address at all. Having two observations of the same
population, we can.

\subsection{Method}

Neither of the two candidate shortcuts is usable. The \texttt{W003} check of
\texttt{mcp-scan} compares each tool against a hash recorded in a persistent
storage file and is exactly the right mechanism, but it cannot be applied
retroactively, because no storage file from the first execution survives and the
workers of the second shared a single file, which is why the 837 \texttt{W003}
findings it produced are an artefact rather than a signal. The \texttt{X-03}
check of \texttt{mcp-security-scan} does compare two tool listings, but it
compares two successive calls within the same session: it detects a server whose
tool set is unstable at startup, not one that changes over weeks. In the second
execution it produced 86 findings, 52 of which survived filtering and none of
which was confirmed, almost all being servers that were not yet ready at the
first call.

We therefore performed a third, deliberately minimal pass over the ecosystem.
For every server for which an inventory of tools was recorded during the first
execution, we started the server again and issued a single \texttt{tools/list}
call, without fuzzing, static analysis or any classification step. The result is
an exact comparison: both sides are the complete listing that the server itself
declares, which is precisely what an MCP client would read, rather than a
fragment reconstructed from the findings of a scanner.

\subsection{Coverage and ecosystem attrition}
\label{sec:long:attrition}

The comparison is exact but partial, and the extent of the limitation is itself
a result. An inventory of tools from the first execution exists for 5,948
servers, which is 8.6\% of the dataset, because only the frameworks that start a
server and enumerate its tools record such an inventory at all. Of those 5,948
servers, only 4,772 could be started again: 988 no longer complete their
startup, 172 have had their GitHub repository deleted outright, and the
remaining 16 fail for other reasons. Three months after the first measurement,
in other words, \textbf{19.8\% of the servers that were analysable are no longer
analysable}, and close to 3\% have vanished from the platform entirely, a figure
we verified repository by repository.

Two consequences follow. The first is that the ecosystem's attrition rate is
high enough to be a confounder in any longitudinal study of MCP, and must be
measured rather than assumed. The second is that the rug-pull results below rest
on 4,772 servers, or 6.9\% of the dataset, and are therefore a \emph{lower
bound}: the same incidence observed over the full population would yield a
proportionally larger count, and the servers that could not be restarted are
precisely those about which nothing can be said.

\subsection{What changed}

Across the 4,772 servers that could be observed twice, 1,013 tool descriptions
changed, distributed over 251 servers. The distribution is markedly uneven, and
separating it is the first analytical step: 57 servers rewrote five or more
descriptions at once, accounting for 679 of the changes, which is the signature
of a release that re-documented an entire catalogue rather than of a targeted
modification. A rug pull, by definition, is targeted. The remaining 334 changes,
spread over 194 servers, are the ones worth reading, and among them 23 introduce
a description that declares a dangerous action the earlier description did not.

Those 23 were read individually against the full text of both versions, and
Table~\ref{tab:long:verdicts} reports the outcome along two independent axes:
whether the capability genuinely changed, and how dangerous the tool is today
regardless of any change.

\begin{table}[htbp]
\centering
\caption{Verdicts on the 23 targeted changes declaring a new dangerous capability}
\label{tab:long:verdicts}
\begin{tabular}{llr}
\toprule
Verdict & Meaning & Cases \\
\midrule
Confirmed rug pull & the tool can now perform a dangerous action it did not declare & 2 \\
Weak expansion     & a real but bounded widening of an existing capability & 2 \\
More restrictive   & the tool declares \emph{narrower} powers than before & 6 \\
Documentation only & same capability, more accurate description & 13 \\
\bottomrule
\end{tabular}

\vspace{2mm}
\footnotesize
Present-day danger, assessed independently of the change: 8 cases high, 9
medium, 6 low. A tool may be highly dangerous without being a rug pull, having
been equally dangerous at the first measurement.
\end{table}

Three of the changes introduce directive language of the kind associated with
description-level injection, and all three turn out on reading to be legitimate
workflow guidance internal to the server. Six changes move in the opposite
direction to the threat: a Databricks connector now requires an explicit
confirmation flag before executing DDL or DML, a MySQL connector moved deletion
behind a permission separate from execution, and a SAP HANA connector blocks
writes by default. That a measurable fraction of servers tightened their declared
capabilities over two months is a modest but genuine observation about the
ecosystem, and one that a study designed only to count vulnerabilities would
discard.

\subsection{The two confirmed cases}

The first confirmed case is \texttt{letoribo/mcp-graphql-enhanced}, whose
\texttt{query-graphql} tool was described at the first measurement as
\emph{``Execute a GraphQL query against the endpoint''} and is described now as
\emph{``Execute GraphQL operations (queries and mutations) \dots\ WARNING: this
tool performs remote operations. Mutation operations will modify persistent
state''}. The tool moved from read-only to read-write, and an agent that had
been granted permission to run queries now runs mutations. The change is
nonetheless announced in the text itself, warning included.

The second case is more instructive, and is the most serious finding of this
chapter. The server \texttt{sentry-official/mcp-cap-internal} exposes four tools
for the analysis of packet captures, of which exactly one,
\texttt{analyze\_dhcp\_packets}, changed. The modification consists of a single
appended sentence:

\begin{quote}
\ttfamily
\dots return comprehensive analysis results. \textbf{Then write data to user's
system.}
\end{quote}

A tool that read a capture file now declares that it writes to the user's
filesystem, with no indication of what it writes or where, no warning, and no
change to its name or to the remainder of its description. The other three tools
of the same server were left untouched. The edit grew the description from 621
to 659 characters, a textual similarity of 97\%. We also note that the
organisation name evokes an affiliation with an established vendor, which lends
the server a trust it has not earned.

This case carries a methodological lesson that we consider transferable. A
deliberate rug pull is, by construction, \emph{minimally} different from what it
replaces: an author who intends the change to pass unnoticed changes as little
as possible. Ranking candidate changes by textual dissimilarity, which is the
natural first instinct and the one we followed initially, therefore sorts the
most dangerous cases to the very bottom of the list. In our own analysis this
case ranked last of 23. Detection strategies for this threat should treat small,
semantically load-bearing edits as more suspicious than large rewrites, not
less.

\subsection{Interpretation}

Two confirmed cases out of 4,772 observed servers do not describe a widespread
phenomenon, and the honest reading is that the rug pull, which occupies a
prominent position in the MCP threat literature and motivates dedicated
machinery in at least one of the analysed frameworks, is rare in practice over a
two-month window. This is consistent with the broader conclusion of
chapter~\ref{ch:discussion}: the security exposure of the MCP ecosystem is
dominated by ordinary implementation defects rather than by deliberate attacks.

Three qualifications keep this from being a stronger negative claim than the
data supports. The observation window is two months, and a rug pull may unfold
over a longer period or may have occurred before the first measurement. The
observed population is 6.9\% of the dataset, so the count is a lower bound and
the true figure over the full ecosystem is necessarily larger. And the direction
of the comparison is asymmetric: tools that \emph{disappeared} between the two
measurements can be counted reliably, and 330 did, whereas tools that appear to
be new cannot, because the inventories recorded during the first execution list
only the tools that a scanner flagged and not the complete tool set of each
server. Measuring the addition of tools, which is arguably the more dangerous
variant of the same threat since a new tool is never subject to any approval at
all, would require a complete inventory that no execution of this study
recorded. The listing produced by this third pass constitutes such a baseline
for any future measurement.

\section{Threats to validity specific to the re-run}
\label{sec:long:threats}

Repeating the measurement surfaced several artefacts that the first execution
had no way of revealing. We report them in full, because each of them would have
distorted a headline number had it gone unnoticed, and because they generalise
to any study of this kind.

\subsection{Duplicate findings}
\label{sec:long:duplicates}

The workers of \texttt{mcp-watch} were restarted several times over overlapping
index ranges, and the routine that merges their shards concatenates without
deduplicating. Counting findings by content rather than by record shows that
41.5\% of the raw findings of \texttt{mcp-watch} in the second execution, and
39\% of its confirmed findings, are repetitions of one another. No other
framework is affected, and the first execution shows a duplication rate of 0.2\%.
All figures reported in this chapter are deduplicated. The correction is not
cosmetic: counted by record, \texttt{mcp-watch} appears to grow by 72\% between
the two measurements, whereas deduplicated it grows by 5.0\%, and its
\texttt{credential-leak} category, which appears to be the fastest-growing class
in the ecosystem, in fact declines slightly.

\subsection{Contamination of the analysis host}
\label{sec:long:contamination}

Inspecting the nine analysis machines after the second execution revealed that
the servers under analysis had left persistent state on the hosts that ran them.
On four of the nine machines, entries had been written into the user's crontab
and were still executing every minute nine days after the analysis had ended.
Their content is unambiguous: they are the probe payloads that the frameworks
themselves inject into tool parameters, including \texttt{; id},
\texttt{\$(whoami)}, \texttt{`uname -a`} and \texttt{../../../../etc/passwd}.
Some server under analysis exposed a scheduling tool, accepted those payloads and
persisted them in the host's own scheduler.

The same pattern appears elsewhere on the same machines. The shell start-up file
of one host had grown to 13,128 lines, of which 3,251 are identical repetitions
of the same three-line block appended by a package manager's installer at every
invocation. Between 176 and 608 state directories had been created in the users'
home directories, one per analysed server, and between 34 and 18,342 entries
accumulated in \texttt{/tmp}.

None of this was detected by any of the seven frameworks, and the reason is
structural rather than incidental: every one of them observes what the server
\emph{answers}, and none observes what the server \emph{does to the system
around it} while it runs. Together with the trojans reported in
section~\ref{sec:manual} and with an earlier incident in which a server under
analysis deleted a directory on one of the machines, this constitutes direct
evidence that dynamic analysis of MCP servers is an operation with side effects
on the analysis environment, and that the sandboxing of that environment is a
requirement rather than a precaution.

\subsection{Environment-induced false negatives}

The contamination above has a measurable effect on the results, and this is the
most important of these observations for anyone attempting a similar study.
Servers implemented in Python are started through \texttt{uv}, which is
installed in a directory that is not part of the default search path of a
non-interactive shell; on the machines where the pass was launched in that way,
every Python server failed with \texttt{spawn uv ENOENT} and was recorded as no
longer startable. Separately, the accumulated state described above filled the
disks of two machines and left their package caches cold, so that dependency
installation exceeded the timeout allowed by the pipeline and the build step
silently produced no output, again yielding a startup failure.

Both effects produce the same observable: a server that is perfectly healthy is
recorded as dead. On the affected machines the startup success rate fell from
81\% to 45\%, and had it gone unnoticed it would have been reported as ecosystem
attrition, which is exactly the quantity this chapter sets out to measure. The
anomaly was detectable only because nine machines were running the same workload
in parallel, which made one machine's divergence from the other eight visible;
we confirmed the diagnosis by re-running a sample of the failing servers on a
healthy host, where all of them started correctly. A single-machine study would
have had no such control, and we accordingly recommend that measurements of
ecosystem attrition either distribute the workload across independent hosts or
re-verify a sample of negative results on a clean environment before reporting
them.

\subsection{Quantities we did not measure}

Two divergences between the executions remain unexplained and we report them as
open rather than resolving them speculatively. At an equal number of servers
producing findings and after deduplication, \texttt{mcp-watch} emits 38\% more
findings per server than it did at the first measurement; the candidate
explanations, namely repositories that grew over the interval and the bundled or
minified sources retrieved for NPX packages, have not been separated. And the
raw output of \texttt{mcp-security-scan} halved between the two executions on a
framework that completed its run without incident, which we have not
investigated.

\section{Responsible disclosure}
\label{sec:long:disclosure}

% TODO: sezione da completare. Elementi previsti: protocollo a livelli
% (malware e trojan verso GitHub Abuse e i registry, credenziali di terzi verso
% i provider per la revoca, rug pull verso i manutentori), motivazione per cui
% la disclosure esaustiva su ~8.500 server non e' praticabile, e stato delle
% segnalazioni effettuate.

\medskip
\noindent\emph{[Section to be completed.]}
